% TOP_PROBLEM: {"id": 30002363, "title": "Lefschetz Classes and the Tate and Hodge Conjectures", "category_id": 5, "category": {"id": 5, "name": "algebraic_geometry", "display_name": "Algebraic Geometry", "description": "Geometric objects defined by polynomial equations.", "slug": "algebraic-geometry", "order_index": 5, "created_at": "2026-07-31T15:26:25.671Z"}, "status": "open"}
% FIRST_POSTED: 2026-10-06
% ATTEMPT_STATUS: solved
% OPENAI_CLAIM: {"schema_version":1,"source_repo":"https://github.com/openai/math","source_commit":"adc7f1241b42e322a6451854ab7e4b4c146bf78a","release_date":"2026-10-06","match":"stronger","resolution":"proved","independently_reviewed":false,"second_pass":"agreed","families":[{"family":"032","title":"Hodge and Kuga–Satake results for all projective K3 surfaces","subject":"Algebraic and complex geometry","manuscripts":[{"title":"The rational Hodge conjecture for CM abelian varieties","pdf_path":"preprints/The-rational-Hodge-conjecture-for-CM-abelian-varieties-September-30-2026/paper.pdf","readme_path":"preprints/The-rational-Hodge-conjecture-for-CM-abelian-varieties-September-30-2026/README.md","date":"2026-09-30","theorem_ref":"thm:cm (sections/01-introduction.tex:17); cc:ghc (sections/07-consequences.tex:154); cc:tate (sections/07-consequences.tex:213); cc:standard (sections/07-consequences.tex:274)","statement_sources":[{"path":"preprints/The-rational-Hodge-conjecture-for-CM-abelian-varieties-September-30-2026/build/sections/01-introduction.tex","line":17,"label":"thm:cm"},{"path":"preprints/The-rational-Hodge-conjecture-for-CM-abelian-varieties-September-30-2026/build/sections/07-consequences.tex","line":154,"label":"cc:ghc"},{"path":"preprints/The-rational-Hodge-conjecture-for-CM-abelian-varieties-September-30-2026/build/sections/07-consequences.tex","line":213,"label":"cc:tate"},{"path":"preprints/The-rational-Hodge-conjecture-for-CM-abelian-varieties-September-30-2026/build/sections/07-consequences.tex","line":274,"label":"cc:standard"}]},{"title":"Weil classes and Hodge classes on abelian powers","pdf_path":"preprints/Weil-classes-and-Hodge-classes-on-abelian-powers-September-30-2026/paper.pdf","readme_path":"preprints/Weil-classes-and-Hodge-classes-on-abelian-powers-September-30-2026/README.md","date":"2026-09-30","theorem_ref":"p3ab:split-sixfolds (paper.tex:79); p3ab:quadratic-five (paper.tex:88)","statement_sources":[{"path":"preprints/Weil-classes-and-Hodge-classes-on-abelian-powers-September-30-2026/build/paper.tex","line":79,"label":"p3ab:split-sixfolds"},{"path":"preprints/Weil-classes-and-Hodge-classes-on-abelian-powers-September-30-2026/build/paper.tex","line":88,"label":"p3ab:quadratic-five"}]}],"lean":null,"reasoning_trace":null}]}
\documentclass[11pt]{article}
\usepackage[margin=25mm]{geometry}
\usepackage{fontspec}
\setmainfont{Latin Modern Roman}
\usepackage{amsmath,amssymb,mathtools}
\usepackage{unicode-math}
\setmathfont{Latin Modern Math}
\usepackage{enumitem}
\usepackage{xurl,hyperref,bookmark}
\hypersetup{colorlinks=true,urlcolor=blue}
\setcounter{secnumdepth}{0}
\setlength{\parindent}{0pt}
\setlength{\parskip}{5pt}
\setlength{\emergencystretch}{3em}
\begin{document}
\section*{Lefschetz Classes and the Tate and Hodge Conjectures}

\subsection{OpenAI's claim}
Manuscript 1 (family 032, 2026-09-30): The manuscripts claim algebraicity of rational Hodge classes on CM abelian varieties, with consequences for generalized Hodge classes and finite-field Tate classes. Manuscript 2 (family 032, 2026-09-30): Algebraicity is claimed for Hodge classes on powers of split Weil sixfolds and of abelian varieties of dimension at most five carrying an imaginary-quadratic action, including special periods.

\subsection{Mathematical statement of the claimed result}
\paragraph{Mathematical claim 1.} A complex abelian variety $A$ is CM if $\operatorname{End}(A)\otimes\mathbb Q$ contains a commutative semisimple algebra of dimension $2\dim A$. For every such $A$ and every $p\ge0$, OpenAI claims
\[
\operatorname{CH}^p(A)_{\mathbb Q}\longrightarrow H^{2p}(A,\mathbb Q)\cap H^{p,p}(A)
\quad\text{is surjective}.
\]
This also applies to arbitrary finite products of CM abelian varieties, including their powers. If $W\subset H^k(A,\mathbb Q)$ is a rational Hodge substructure and $W(r)$ is effective (all Hodge bidegrees are nonnegative), it claims a closed algebraic $Z\subset A$ of codimension at least $r$ with $W\subset\ker[H^k(A,\mathbb Q)\to H^k(A\setminus Z,\mathbb Q)]$.
The manuscript additionally asserts that, for every abelian variety $A/\mathbb F_q$, every $r\ge0$ and prime $\ell\ne\operatorname{char}\mathbb F_q$,
\[
\operatorname{CH}^r(A)\otimes\mathbb Q_\ell\longrightarrow
H^{2r}_{\mathrm{\acute et}}(A_{\overline{\mathbb F}_q},\mathbb Q_\ell(r))^{\operatorname{Gal}(\overline{\mathbb F}_q/\mathbb F_q)}
\quad\text{is surjective}.
\]
Numerical and $\ell$-adic homological equivalence are claimed to agree on rational cycles. Over any algebraically closed field, for an abelian $g$-fold and rational ample class $h$, the form $(-1)^r\deg(xyh^{g-2r})$ is claimed positive definite on the real extension of $\ker[h^{g-2r+1}:A^r_{\mathrm{num}}\to A^{g-r+1}_{\mathrm{num}}]$ for $0\le r\le g/2$.

\paragraph{Mathematical claim 2.} Let $E=\mathbb Q(u)$ be imaginary quadratic with $u^2=-d$, $d>0$, and let $A$ be a complex abelian variety with a unital $E$-action. OpenAI claims rational cycle-class surjectivity
\[
\operatorname{CH}^p(A^m)_{\mathbb Q}\longrightarrow H^{2p}(A^m,\mathbb Q)\cap H^{p,p}(A^m)
\quad(m\ge1,\ p\ge0)
\]
in either of these two cases: $\dim A\le5$; or $\dim A=6$ and the rational Hermitian form associated to a compatible polarization is hyperbolic with signature $(3,3)$. In the second case, with rational Riemann form $e$, this means $h_E(x,y)=e(x,uy)+u e(x,y)$ on $H_1(A,\mathbb Q)$ has a totally isotropic three-dimensional $E$-subspace. Compatibility means $e(ax,y)=e(x,\bar a y)$ for all $a\in E$.

\subsection{Outline of the argument}
Manuscript 1 (family 032, 2026-09-30): The argument realizes four-factor balanced tensors using unitary theta constructions. Elementary switches between equal-sum lists, followed by algebraic polarization contractions and descent of coefficients, build arbitrary balanced CM tensors from these elementary cycle classes.

Manuscript 2 (family 032, 2026-09-30): A coherent seed and its Fourier--Mukai transform produce the split-sixfold Weil plane through an obstruction calculation and deformation. A determinant description of the actual Hodge group, combined with the CM result, extends these algebraic generators to powers and to the stated low-dimensional cases.

\subsection{Scope relative to this problem}
The CM theorem claims algebraic cycle representatives for all rational Hodge classes on complex CM abelian varieties, including their finite products and powers. Its Tate corollary claims cycle-class surjectivity for all abelian varieties defined over finite fields, hence for powers of CM reductions. Both conclusions hold without the recorded antecedent concerning generation by Lefschetz classes. Removing that sufficient-condition hypothesis gives a stronger assertion.

\subsection{Formal verification}
Family 032 has no Lean scope page in this snapshot; no formalization of these manuscript results is listed.

This site has not executed the Lean files or checked their dependencies. These descriptions report the scope of the linked declarations, rather than a new verification of the proofs.

\subsection{Status caveat}
OpenAI attributes these results to an internal model. Its release notice cautions that results lacking full formalization can contain mistakes. This site has not independently assessed the proofs; a solved label records the scope of the reported claim.

\subsection{Sources}
\begin{itemize}
\item OpenAI, \emph{The rational Hodge conjecture for CM abelian varieties}, 2026-09-30; citation key \texttt{OAI:The-rational-Hodge-conjecture-for-CM-abelian-varieties-September-30-2026}. \href{https://github.com/openai/math/blob/adc7f1241b42e322a6451854ab7e4b4c146bf78a/preprints/The-rational-Hodge-conjecture-for-CM-abelian-varieties-September-30-2026/paper.pdf}{PDF (GitHub)}; \href{https://github.com/openai/math/raw/adc7f1241b42e322a6451854ab7e4b4c146bf78a/preprints/The-rational-Hodge-conjecture-for-CM-abelian-varieties-September-30-2026/paper.pdf}{PDF (direct)}; \href{https://github.com/openai/math/blob/adc7f1241b42e322a6451854ab7e4b4c146bf78a/preprints/The-rational-Hodge-conjecture-for-CM-abelian-varieties-September-30-2026/README.md}{Manuscript page and citation}; \href{https://github.com/openai/math/blob/main/preprints/The-rational-Hodge-conjecture-for-CM-abelian-varieties-September-30-2026/paper.pdf}{Latest version}; Statement source: \href{https://github.com/openai/math/blob/adc7f1241b42e322a6451854ab7e4b4c146bf78a/preprints/The-rational-Hodge-conjecture-for-CM-abelian-varieties-September-30-2026/build/sections/01-introduction.tex\#L17}{thm:cm}; Statement source: \href{https://github.com/openai/math/blob/adc7f1241b42e322a6451854ab7e4b4c146bf78a/preprints/The-rational-Hodge-conjecture-for-CM-abelian-varieties-September-30-2026/build/sections/07-consequences.tex\#L154}{cc:ghc}; Statement source: \href{https://github.com/openai/math/blob/adc7f1241b42e322a6451854ab7e4b4c146bf78a/preprints/The-rational-Hodge-conjecture-for-CM-abelian-varieties-September-30-2026/build/sections/07-consequences.tex\#L213}{cc:tate}; Statement source: \href{https://github.com/openai/math/blob/adc7f1241b42e322a6451854ab7e4b4c146bf78a/preprints/The-rational-Hodge-conjecture-for-CM-abelian-varieties-September-30-2026/build/sections/07-consequences.tex\#L274}{cc:standard}.
\item OpenAI, \emph{Weil classes and Hodge classes on abelian powers}, 2026-09-30; citation key \texttt{OAI:Weil-classes-and-Hodge-classes-on-abelian-powers-September-30-2026}. \href{https://github.com/openai/math/blob/adc7f1241b42e322a6451854ab7e4b4c146bf78a/preprints/Weil-classes-and-Hodge-classes-on-abelian-powers-September-30-2026/paper.pdf}{PDF (GitHub)}; \href{https://github.com/openai/math/raw/adc7f1241b42e322a6451854ab7e4b4c146bf78a/preprints/Weil-classes-and-Hodge-classes-on-abelian-powers-September-30-2026/paper.pdf}{PDF (direct)}; \href{https://github.com/openai/math/blob/adc7f1241b42e322a6451854ab7e4b4c146bf78a/preprints/Weil-classes-and-Hodge-classes-on-abelian-powers-September-30-2026/README.md}{Manuscript page and citation}; \href{https://github.com/openai/math/blob/main/preprints/Weil-classes-and-Hodge-classes-on-abelian-powers-September-30-2026/paper.pdf}{Latest version}; Statement source: \href{https://github.com/openai/math/blob/adc7f1241b42e322a6451854ab7e4b4c146bf78a/preprints/Weil-classes-and-Hodge-classes-on-abelian-powers-September-30-2026/build/paper.tex\#L79}{p3ab:split-sixfolds}; Statement source: \href{https://github.com/openai/math/blob/adc7f1241b42e322a6451854ab7e4b4c146bf78a/preprints/Weil-classes-and-Hodge-classes-on-abelian-powers-September-30-2026/build/paper.tex\#L88}{p3ab:quadratic-five}.
\item \href{https://github.com/openai/math/blob/adc7f1241b42e322a6451854ab7e4b4c146bf78a/CONTENTS.md}{OpenAI catalogue entry}.
\end{itemize}
\end{document}
